\section{序論}

中心の各閉点で量子化を理解できても、その局所表示を境界全体で貼り合わせられるとは限らない。
本稿では全中心と単純加群の合流を求め、さらに中心の解消・完備化に残るBrauer障害を計算する。
特に最深点では、中心ファイバーの単純加群が一次元一つになる一方、
完備局所分数体上のBrauer指数は量子化のPI次数と同じ最大値を保つ。

重み付き射影空間の非可換座標環には、射影圏の変形だけでは見えない
中心と表現の幾何がある。本稿では、非線形な局所冪零導分から作る
四次元CGG量子化族について、この幾何を正規元の局所化により具体的に求める。
特に重み $m=2j+1\ge3$ の $r=1$ 共鳴のCalabi--Yau代表では、
中心に $D_{m+4}$ 型曲面が現れ、その原点の上に非自明なprimitive spectrumが生じる。

まず正共鳴の非PI部分について、主要な結論を三点にまとめる。
第一に、正規元 $x$ を反転した環の中心は $A_{m+3}$ 型であり、
元の $D_{m+4}$ 型中心から一つの例外曲線 $\mathbb P^1$ を取り出す
crepantな部分解消の一つのchartになる。
第二に、中心の原点以外の全てのcentral reductionは
$\mathcal D(\mathbb G_m)$ に同型である。中心の分数体上でも根の付加を必要としない。
第三に、全てのprimitive idealを分類し、中心原点上では
\[
 \operatorname{Prim}_0(\mathcal S)
 =\{\mathfrak p_\lambda:\lambda\in\mathbb P^1(\mathbb C)\}
 \ \sqcup\
 \{\mathfrak n_{z_0,w_0}:(z_0,w_0)\in\mathbb A^2(\mathbb C)\}
\]
という表示を得る。
後者は全ての有限次元単純加群から来る文字表現の族である。
本稿が分類するのは単純加群の零化イデアルであり、
同じ零化イデアルを持つ無限次元単純加群まで同型類ごとに列挙するものではない。

原点上の閉包関係も決定できる。
$\mathfrak p_\lambda$ の一点閉包には、その点自身に加えて
\[
 \{\mathfrak n_{z_0,0}:z_0\in\mathbb C\}
\]
という共通の直線が含まれる。
一方、$\mathfrak p_\lambda$ の族の任意の無限部分集合は、原点上のprimitive spectrum全体に稠密である。
従って、$w_0\ne0$ の文字表現は一つの $\mathfrak p_\lambda$ の特殊化ではないが、
無限個の $\lambda$ を合わせると閉包に現れる。
この閉包公式によって、原点fiber内の任意の部分集合の位相を記述できる。

さらに全ての正共鳴で、量子素スペクトルとPoisson素スペクトルを同じ部分空間として同定する。
原始スペクトルは代数的な葉の空間であり、葉の飽和を取ってからZariski閉包を取る公式が、
非原点を動く中心指標の集積も含めて全域的な位相を記述する。
$r=1$ の原点の根基については、正確な冪零指数 $(m+7)/2$ も得る。
このスキーム論的な厚みは、原始スペクトルの位相とは別の情報である。

\subsection{方法と先行研究}

CGG積を使ったAS正則代数の構成と、正規元局所化によるスペクトルの研究には、
Lecoutre--Sierra\cite{LecoutreSierra2019}の先行研究がある。
彼らの線形な局所冪零導分による族と、本稿の非線形共鳴族は異なるが、
Ore拡大と可解Lie作用を使う方法を共有する。
本稿では局所化中心の指数半群、明示的な逆座標変換、
可換saturationと星積イデアルの一致を組み合わせる。

局所化後の単純加群を元の環へ制限しても、単純性は一般に保存されない。
この点は、可算次元複素整域のJacobson根基が零であることを使う
正規局所化の補題によって処理する。
境界の包含関係には、局所化した曲線の有限整拡大を明示してlying-overを適用する。
これにより、$x\ne0$ の計算だけから境界を推測せずに済む。

全primitive idealの表示から、Dixmier--Moeglin同値も導く。
これは本例のrational primeが全て明示したリストに入り、
正規局所化により局所閉になることによる。
一般のNullstellensatzの含意についてはBell--Wang--Yee\cite{BellWangYee2019}を参照する。

本稿の量子化族のAS正則性と全中心の生成性は、公表稿IX\cite{VKPoisson34}を入力とする。
射影圏の導来分類と変形環は、公表稿X\cite{VKCategories36}にまとめた。
ここでは中心の局所化、部分解消、central reduction、全正共鳴のスペクトル対応と原始イデアルの証明を
独立した一つの論証として記述する。

\subsection{範囲について}

この正共鳴の非PI部分で扱う代数は中心上有限ではなく、通常のAzumaya代数ではない。
中心がDu Val曲面であることからNCCRや非可換解消を自動的に得るとも主張しない。
部分解消の例外曲線と $\mathfrak p_\lambda$ のパラメータは共に $\mathbb P^1$ となり、
有限座標は局所中心の $c=qJ/x$ によって一致する。
しかし部分解消全体の上に代数の層を構成することは、ここでのprimitive idealの分類とは別の課題である。

\subsection{端点のPIファイバーとMcKay模型}

偶数端点の全変形を実現するAS正則CY族では、一般のパラメータは非PIである一方、
一の冪根では任意に高いPI次数を持つファイバーが現れる。
その中心の $x^\ell\ne0$ という開部分は、混合項を消去した一般化Weyl代数で計算できる。
位数 $\ell$ が奇数なら中心は滑らかで全体がAzumaya、
偶数 $2d$ なら $A_1$ 型の特異部分が正確に非Azumaya部分となる。
特異中心ファイバーは、二頂点と各向き二本の矢を持つ平方零代数の $d$ 次行列環である。

この表現の構造を、中心の有限エタール被覆上での滑らかな四次元環と
$\mu_2$ 作用へ持ち上げる。これにより、ファイバーだけでなく完備近傍全体が
McKay skew group模型で説明できる。
この完成局所表示が元の中心全体での自己準同型環表示を意味するかを調べると、
明確な障害がある。一般点の代数は
\[
 D=(\Xi,U)_t,\qquad
 x^\ell=\Xi,\quad u^\ell=U,\quad ux=txu
\]
という巡回斜体で、そのBrauer周期と指数はいずれも $\ell$ である。
従って、元の中心上で有限反射的加群の自己準同型環になることも、そのような環と
中心上でMorita同値になることもない。
完成局所的なNCCRと大域的な自己準同型環の相違を、この周期が検出する。

また特異部分 $\Sigma\simeq\mathbb G_{m,Q}\times\mathbb A^1_J$ 上の全代数は、
$L=\mathbb C[Q^{\pm1},J,\kappa]/(\kappa^2+tQ/\beta)$ による
\[
 A_\Sigma\simeq M_d\bigl(L\ltimes(L_\iota\oplus L_\iota)\bigr),
 \qquad\iota(\kappa)=-\kappa
\]
である。二つの単純加群はこの連結二重被覆上で分離し、$Q$ の一周によって交換される。
式は特異部分への制限全体について成り立ち、閉点ごとの模型だけに留まらない。

\subsection{混合項を回収する境界と巡回被覆}

関数斜体と $x^\ell\ne0$ の中心からは消えていた $H$ を検出するには、
別の中心開部分 $U=u^\ell\ne0$ が有効である。
ここで $H=h(Q,u)+xG$ と分けて $W=w+G$ と置くと、
\[
 xW=J-h(Q,u),\qquad Wx=J-h(Q,tu)
\]
という一般化Weyl表示となる。その全中心は
\[
 \mathbb C[Q,J,U^{\pm1},\Xi,\Omega]/(\Xi\Omega-P_h),
 \qquad P_h=\prod_{i=0}^{\ell-1}(J-h(Q,t^iu))
\]
であり、$\Omega=W^\ell$ である。
従って $H$ は中心の境界への延長を変える。
$Q,U,J$ を固定して $P_h$ を知れば、$h$ は $u\mapsto t^ku$ まで決まる。
重み四・位数五では、$H=u^5$ と $H=u^4v$ が同じ関数斜体と同じ
$x^\ell\ne0$ の開部分を持ちながら、$U\ne0$ の中心の特異部分の次元は異なる。

境界ファイバーを決めるのは $J=h(Q,t^i\lambda)$ が成立する巡回軌道上の位置である。
成立する辺で巡回箙を切り、残る鎖の長さを $d_1,\ldots,d_r$ とすると、
単純加群の次元は正確に $d_1,\ldots,d_r$ であり、根基の冪零指数は $r$ となる。
さらに $h$ の巡回安定化群が位数 $g$ なら、衝突を除いた正規化境界上の
全代数と、連結巡回 $g$ 重被覆上の単純表現族を明示できる。
この部分でもAzumaya部分は中心の滑らかな部分と一致する。

一般化Weyl代数の重み加群と巡回軌道の理論はDrozd--Guzner--Ovsienko
\cite{DGOBoundary39}に帰する。
一の冪根での量子代数の中心・Azumaya部分の具体的研究として
Cooney--Ganev--Jordan\cite{CGJBoundary39}も参照する。
本稿では、この枠組みを混合Ore族の境界へ適用し、全ファイバーの行列表示、
標識付きノルム復元、および元の中心上への降下を明示する。
巡回代数のノルム分裂条件\cite{MinacSchultzSwallowDescent39}から得るBrauer障害と、
境界の有限被覆が記録する情報を、次の形式近傍の計算によって結び付ける。

\subsection{全中心による区別と最深点に残る指数}

混合項が $H=f(Q,u)+xG$ と書ける場合には、$W=w+G$ と直し、
$\Omega=W^\ell$ を加えることで全中心を得る。
巡回ノルムを $N_f(Q,U,J)=\prod_i(J-f(Q,t^iu))$ とすると、
\[
 Z(S_H)=\frac{\mathbb C[Q,J,\Xi,\Omega,U,V]}
 {(UV-P_\ell(Q,\Xi),\ \Xi\Omega-N_f(Q,U,J))}
\]
である。正規な二つの開部分だけでなく、残る余次元二の部分も処理することで、
この表示が局所化前の全中心であることを証明する。

重み四、$\ell=5$、$\beta=1$ で $f=0$ と $f=Qu^3$ を比べると、
第二の方程式はそれぞれ $\Xi\Omega=J^5$ と
$\Xi\Omega=J^5-Q^5U^3$ となる。
前者の特異部分は二次元成分を持つが、後者の被約特異部分は三本の座標軸である。
従って標識を忘れた全中心も非同型であり、二つのアフィン代数は
Morita同値でも $\mathbb C$ 線形dg導来Morita同値でもない。
両者の関数斜体と $\Xi\ne0$ の中心開部分は同じである。
射影圏の同値を排除するには、さらに次数付きの不変量が必要となる。

混合項が $H=H_c+xG$、$H_c\in Z(B)$ と書ける場合には、$w\mapsto w+G$ により
$H=0$ の部分族へ移る。その全中心は
\[
 Z_0=\frac{\mathbb C[Q,J,\Xi,\Omega,U,V]}
 {(\Xi\Omega-J^\ell,\ UV-P_\ell(Q,\Xi))}
\]
であり、$\ell$ の偶奇により
\[
 P_\ell=-Q^\ell-\beta^\ell\Xi^2\quad(\ell\text{ 奇数}),\qquad
 P_{2d}=(Q^d+(-1)^d\beta^d\Xi)^2
\]
となる。これは四次元の正規Gorenstein完全交叉である。
全ての閉中心ファイバーを分類し、Azumaya部分が中心の滑らかな部分と一致すると示す。
最深点では一つの一次元単純加群しか残らないが、ファイバーの全次元は
奇数時 $(7\ell^2-6\ell+1)/2$、偶数時 $(7\ell^2-6\ell+2)/2$ へ増大する。
根基の冪零指数は正確に $2\ell-1$ である。

この原点で中心を完備化しても、一般点の巡回代数の周期と指数は $\ell$ のままである。
証明では $\operatorname{wt}(Q,J,\Xi,U,V)=1$、$\operatorname{wt}(\Omega)=\ell-1$ の
付値を作り、剰余体の独立変数 $U/\Xi$ がノルムであることを妨げる。
従ってこの完成代数は、中心上の反射的加群の自己準同型環とも、その中心線形Morita型ともならない。
$\Xi\ne0$ の完成局所McKay模型を、最深点にまで外挿することはできない。
$f$ が $Q$ の倍数である場合にも、最深点のファイバーと完備局所指数は同じである。
従って上記の全中心が非同型な二例は、これらの最深点の数値も共有している。

\subsection{例外因子の留数としての巡回モノドロミー}

一般の混合項についても、$U\ne0$ の一般境界では全形式近傍を記述できる。
$h$ の安定化群位数を $g$、$n=\ell/g$、$z=u^g$ とすると、
正規化境界の関数体 $F=\mathbb C(Q,z)$ 上の横断中心は
\[
 R_F=F[[\xi,\upsilon,s]]/(\xi\upsilon-s^g)
\]
となる。完成代数は巡回拡大 $L=F[\kappa]/(\kappa^g-z)$ を係数とする角の $n$ 次行列環であり、
$L$ への拡大後には $L[[a,b]]\#\mu_g$ という標準McKay模型になる。
特に一般境界の複素閉点完成はNCCRを持つ。

しかし一般境界点の最小可換解消では、第 $i$ 例外曲線に
\[
 \partial_{E_i}[D]=i\chi,\qquad
 \operatorname{ord}\partial_{E_i}[D]=g/\gcd(g,i)
\]
というBrauer留数が残る。ここで $\chi$ は $L/F$ に対応する巡回指標である。
この式は標準的な巡回代数の留数公式\cite[Proposition~5.8]{HarariBrauer40}の本族への適用である。
全ての留数を消す最小有限係数被覆は $L/F$ であり、
単純加群の巡回モノドロミーを解く被覆と一致する。
これにより、閉点完成のNCCR、一般境界での分岐、最深点に残る指数を一つの計算体系で比較できる。
